\documentclass[book.tex]{subfiles}
\begin{document}

\chapter{Hamiltonian Neural Networks}
\label{chap:hamiltonian_nn}
\noindent\rule{11cm}{0.4pt} \\
\textbf{Prerequisites:} \autoref{chap:hamiltonians_mechanics}, \autoref{chap:neural_ode} \\
\textbf{Difficulty Level:} ** \\
\noindent\rule{11cm}{0.4pt}
\vspace{5mm}

In chapter~\ref{chap:neural_ode}, we learned about neural ordinary differential equations which provide tools to learn the evolution equations of dynamical systems. One limitation of these methods is that they cannot learn invariances or richer structure present within a physical system easily. As a result, such systems may be data inefficient. In cases where we know more about the conserved quantities in a system, we may be able to learn a richer representation.


\section{Parameteric Hamiltonians}

Classically, the function $\mathcal{H}$ is written down analytically by studying the physics of the system (as we have done in earlier chapters). The challenge to this approach is that for complex systems it may not be possible to construct an analytical Hamiltonian. An alternative approach constructs a parameteric Hamiltonian, where $\mathcal{H}$ is parameterized by learnable weights $\theta$. We can then use these weights to construct a natural loss function

\begin{align}
    \mathcal{L}_{\textrm{Hamiltonian}} &= \left \| \frac{\partial \mathcal{H}_{\theta}}{\partial p} - \frac{\partial q}{\partial t}\right \|_2
    +
    \left \|\frac{\partial \mathcal{H}_{\theta}}{\partial q} + \frac{\partial p}{\partial t}\right \|_2
\end{align}

\section{A Mass-Spring System}

Let's study a few simple Hamiltonian systems. We start with a simple ideal mass-spring system as shown in Figure~\ref{fig:ideal_mass_spring}. The Hamiltonian for this system is given by

\begin{align}
    \mathcal{H}_\theta &= \frac{1}{2}kq^2 + \frac{p^2}{2m}
\end{align}
Here $q$ and $p$ are the canonical coordinates and represent the displace of the mass and its momentum respectively. The parameters in this system are given by
\begin{align}
    \theta &= (k, m)
\end{align}
We can write down the explicit form of the loss in this case as well
\begin{align}
    \mathcal{L}_{\textrm{Hamiltonian}} &= \left \| \frac{\partial \mathcal{H}_{\theta}}{\partial p} - \frac{\partial q}{\partial t}\right \|_2
    +
    \left \|\frac{\partial \mathcal{H}_{\theta}}{\partial q} + \frac{\partial p}{\partial t}\right \|_2 \\
    &= \left \|\frac{p}{m} - \dot{q} \right \|_2 + \| kq + \dot{p} \|_2 \\
    &= \left \|\frac{p}{m} - \dot{q} \right \|_2 + \| kq + m\ddot{q} \|_2 
\end{align}
%(\textcolor{red}{TODO: Is there a sign error?}).
\begin{marginfigure}
    \centering
    % \includegraphics{figures/Differentiable Physics/Mechanics_ML/Hamiltonian_NN/ideal_mass_spring.png}
     \includegraphics{figures/Differentiable Physics/Mechanics_ML/Hamiltonian_NN/Mass Spring System (1).png}
    \caption{Ideal mass spring system.}
    \label{fig:ideal_mass_spring}
\end{marginfigure}

For a synthetic dataset, we assume that the true dynamics are obscured by Gaussian noise and our task is to learn a model of the dynamics from the noisy dataset.

\section*{Exercises}
\begin{enumerate}
    \item Write down the Hamiltonian of a system of two coupled springs (see Figure~\ref{fig:coupled_springs}). Write down the explicit loss function for this system.
    
    \begin{marginfigure}
        \centering
        % \includegraphics{figures/Differentiable Physics/Mechanics_ML/Hamiltonian_NN/coupled_springs.png}
         \includegraphics[width=\linewidth]{figures/Differentiable Physics/Mechanics_ML/Hamiltonian_NN/Coupled springs.png}
        \caption{A system of two coupled springs}
        \label{fig:coupled_springs}
    \end{marginfigure}
\end{enumerate}

\end{document}